\section{Análisis del Dataset y Dimensionalidad}
\label{sec:eda}

\subsection{Descripción del Dataset}
El conjunto de datos \textit{Brain Tumor Classification MRI}
contiene \textbf{3\,264 imágenes} MRI cerebrales distribuidas en
cuatro clases, particionadas en conjuntos de entrenamiento y
prueba por los autores originales. La Tabla~\ref{tab:dist}
resume su distribución.

\begin{table}[H]
\centering
\caption{Distribución de imágenes por clase y partición.}
\label{tab:dist}
\begin{tabular}{lrrr}
\toprule
\textbf{Clase} & \textbf{Train} & \textbf{Test} & \textbf{Total} \\
\midrule
glioma\_tumor      & 826 & 100 & 926 \\
meningioma\_tumor  & 822 & 115 & 937 \\
no\_tumor          & 395 & 105 & 500 \\
pituitary\_tumor   & 827 &  74 & 901 \\
\midrule
\textbf{Total}     & \textbf{2\,870} & \textbf{394} & \textbf{3\,264} \\
\bottomrule
\end{tabular}
\end{table}

El conjunto de entrenamiento concentra el \textbf{87{,}9\,\%} de las
imágenes (2\,870), mientras que el conjunto de prueba representa
el \textbf{12{,}1\,\%} restante (394 imágenes). Esta partición
oficial se respeta como \textit{holdout} independiente para la
evaluación final, y el conjunto de entrenamiento se sub-divide
internamente para validación (Sección~\ref{sec:proto}).

\subsection{Balance de Clases}
El análisis de balance evidencia que la clase \textit{no\_tumor}
es la minoritaria, con aproximadamente \textbf{395} imágenes en
entrenamiento, frente a $\sim$\textbf{825} para cada una de las
clases tumorales. Esta proporción asimétrica ($\sim$1:2) puede
inducir sesgos durante el aprendizaje, llevando al modelo a
sub-representar la clase minoritaria.

Como estrategia de compensación se adopta el uso de pesos por
clase (\textit{class\_weight}) en la función de pérdida
\textit{Cross-Entropy}, con ponderaciones inversamente
proporcionales a la frecuencia relativa de cada clase en el
conjunto de entrenamiento. Esto incrementa la magnitud del
gradiente para errores cometidos sobre la clase minoritaria sin
alterar la distribución de muestras observadas. Las
Figuras~\ref{fig:balance} y \ref{fig:pie} ilustran gráficamente
esta distribución.

\begin{figure}[H]
\centering
\includegraphics[width=\columnwidth]{balance_clases}
\caption{Distribución absoluta de imágenes por clase y partición.}
\label{fig:balance}
\end{figure}

\begin{figure}[H]
\centering
\includegraphics[width=0.75\columnwidth]{pie_balance}
\caption{Proporción relativa de las clases en el conjunto de
entrenamiento.}
\label{fig:pie}
\end{figure}

\subsection{Análisis de Dimensionalidad}
Se inspeccionó una muestra de 80 imágenes por clase para
caracterizar la variabilidad en resolución espacial. El dataset
presenta resoluciones altamente heterogéneas: las dimensiones
máximas observadas alcanzan \textbf{1024$\times$768\,px} y las
mínimas \textbf{59$\times$80\,px}. Además, la relación de aspecto
varía considerablemente entre estudios, debido a las diferencias
entre equipos de adquisición y protocolos clínicos.

Dada esta heterogeneidad, se establece el redimensionamiento
uniforme a \textbf{224$\times$224\,px} como paso obligatorio del
pipeline. Este tamaño constituye un compromiso entre detalle
anatómico preservado y costo computacional, y es además el formato
canónico de las arquitecturas CNN modernas, lo que facilita
comparaciones posteriores. Las Figuras~\ref{fig:hist-res} y
\ref{fig:aspect} muestran respectivamente el histograma de
resoluciones y la distribución de relaciones de aspecto.

\begin{figure}[H]
\centering
\includegraphics[width=\columnwidth]{histograma_resolucion}
\caption{Histograma de resoluciones (ancho $\times$ alto) en el
dataset.}
\label{fig:hist-res}
\end{figure}

\begin{figure}[H]
\centering
\includegraphics[width=\columnwidth]{aspecto_ratio}
\caption{Distribución de relaciones de aspecto observadas por
clase.}
\label{fig:aspect}
\end{figure}

\subsection{Análisis Estadístico de Intensidades}
Se analizaron las distribuciones de intensidad de píxeles (escala
de grises) y canales RGB por clase, con el fin de identificar
diferencias fotométricas discriminativas. La
Tabla~\ref{tab:stats} resume los estadísticos descriptivos por
clase.

\begin{table}[H]
\centering
\caption{Estadísticas de intensidad (escala 0--255) por clase.}
\label{tab:stats}
\begin{tabular}{lrrrrr}
\toprule
\textbf{Clase} & \textbf{Media} & \textbf{Std} & \textbf{Min} & \textbf{Max} & \textbf{Mediana} \\
\midrule
glioma\_tumor     & 36{,}23 & 41{,}43 & 0 & 255 & 14 \\
meningioma\_tumor & 44{,}56 & 47{,}77 & 0 & 255 & 28 \\
no\_tumor         & 52{,}21 & 59{,}63 & 0 & 255 & 22 \\
pituitary\_tumor  & 49{,}50 & 42{,}94 & 0 & 255 & 50 \\
\bottomrule
\end{tabular}
\end{table}

Se observa que los \textit{gliomas} presentan los valores más
bajos de media e intensidad mediana, lo que los caracteriza como
tejidos relativamente \textbf{hipointensos} en la modalidad
analizada. Los tumores \textit{pituitarios}, en contraste,
presentan la mediana más alta (50) y la menor dispersión relativa,
indicando intensidades más concentradas y homogéneas. La clase
\textit{no\_tumor} muestra la mayor desviación estándar (59{,}63),
consistente con la variedad anatómica del cerebro sano sin
focalización lesional. Estas diferencias fotométricas constituyen
señales útiles que una CNN puede explotar a nivel de filtros
convolucionales de bajo nivel.

Para estabilizar el entrenamiento se aplicará normalización
\textit{Z-score} sobre los tres canales utilizando los
estadísticos canónicos de ImageNet ($\mu = [0{,}485, 0{,}456,
0{,}406]$, $\sigma = [0{,}229, 0{,}224, 0{,}225]$). Las
Figuras~\ref{fig:box} y \ref{fig:rgb} complementan visualmente
este análisis.

\begin{figure}[H]
\centering
\includegraphics[width=\columnwidth]{boxplot_intensidades}
\caption{Diagramas de caja de intensidades de píxel por clase.}
\label{fig:box}
\end{figure}

\begin{figure}[H]
\centering
\includegraphics[width=\columnwidth]{canales_rgb}
\caption{Distribución de intensidades por canal RGB y por clase.}
\label{fig:rgb}
\end{figure}

\subsection{Imagen Promedio y Muestras por Clase}
La Figura~\ref{fig:promedio} presenta la imagen promedio por
clase, calculada como la media píxel a píxel de 100 imágenes
redimensionadas a 224$\times$224\,px. Los gliomas exhiben
regiones \textbf{hipointensas difusas}, de bordes poco definidos,
consistente con su naturaleza infiltrativa. Los tumores
\textit{pituitarios} tienden a concentrarse en la zona central
de la imagen, reflejando la ubicación anatómica de la glándula
hipófisis. Los meningiomas muestran patrones intermedios con
tendencia a localizarse en la periferia. La clase
\textit{no\_tumor}, al promediar tejidos cerebrales sanos con
distintas orientaciones, produce una imagen más difusa y
centralmente simétrica.

\begin{figure}[H]
\centering
\includegraphics[width=\columnwidth]{imagen_promedio}
\caption{Imagen promedio por clase (100 imágenes,
224$\times$224\,px).}
\label{fig:promedio}
\end{figure}

La Figura~\ref{fig:muestras} ilustra cinco ejemplos
representativos por clase. Se evidencia alta variabilidad
intra-clase en orientación, escala y contraste, lo que justifica
el uso de \textit{data augmentation} durante el entrenamiento
para mejorar la capacidad de generalización del modelo.

\begin{figure}[H]
\centering
\includegraphics[width=\columnwidth]{muestras_por_clase}
\caption{Muestras representativas (5 por clase) del conjunto de
entrenamiento.}
\label{fig:muestras}
\end{figure}

\subsection{Plan de Preprocesamiento}
A partir de los hallazgos anteriores, se define el pipeline de
preprocesamiento resumido en la Tabla~\ref{tab:pipe}.

\begin{table}[H]
\centering
\caption{Pipeline de preprocesamiento propuesto.}
\label{tab:pipe}
\begin{tabularx}{\columnwidth}{clX}
\toprule
\textbf{\#} & \textbf{Paso} & \textbf{Justificación} \\
\midrule
1 & Resize 224$\times$224\,px & Resoluciones heterogéneas requieren entrada uniforme. \\
2 & Conversión a RGB & Algunas imágenes están en escala de grises; se unifica el formato. \\
3 & Z-score ($\mu$ ImageNet) & Estabiliza el gradiente y mejora la convergencia. \\
4 & \textit{Data augmentation} & Rotación $\pm 15^{\circ}$, flip horizontal y variación de brillo: compensa la variabilidad intra-clase. \\
5 & \textit{class\_weight} & Compensa el desbalance de la clase \textit{no\_tumor}. \\
6 & Split 80/20 \textit{Training} & Separa validación interna sin contaminar el \textit{holdout} oficial. \\
\bottomrule
\end{tabularx}
\end{table}
